% Smoke test — exercises every macro in pyportfolios-note.cls once.
% Not a deliverable; used to catch class errors before building a real note.
\documentclass{pyportfolios-note}

\noteshorttitle{Risk parity vs. inverse volatility}

\begin{document}
\thispagestyle{notecover}

\notemasthead{Quantitative Insights}{Research note}{Q2 2026}{Out-of-sample period: 2017–2023}

\notetitle{Risk parity, or just\\inverse volatility?}

\notedeck{When seven correlated mega-cap names sit at the heart of every passive
equity book, two parameter-light allocators land within a hundredth of a Sharpe
of each other — but only one of them is actually doing what it claims.}

\begin{multicols}{2}
\notelede{Risk parity has acquired the status of a portfolio-construction default.
The pitch is straightforward: weight assets so that each contributes equally to
total portfolio variance.}

This note compares the two methods on a deliberately narrow universe: seven US
mega-cap technology stocks (AAPL, MSFT, GOOG, AMZN, META, TSLA, NVDA).

\notesection{The two methods}
\notesubsection{Inverse volatility}
The simplest of the parameter-light allocators. Each asset's weight is set
proportional to the reciprocal of its standalone volatility.

\bottomline{On out-of-sample tech, the two allocators land within hundredths of a
Sharpe. The performance gap is economically negligible, but the
variance-attribution profile is meaningfully different.}
\end{multicols}

\newpage
\notesection{Out-of-sample results}

\begin{kpirow}
\pqkpi{RISK PARITY}{1.12}{Annualised Sharpe}\hfill
\pqkpi{INVERSE VOL}{1.10}{Annualised Sharpe}\hfill
\pqkpi{SHARPE GAP}{0.02}{Risk parity advantage}\hfill
\pqkpi{MAX DRAWDOWN}{56.6\%}{vs. 55.7\% for IV}
\end{kpirow}

\begin{metrictable}{Metric}{Risk parity}{Inverse volatility}{Difference}
\bandrow{RETURN}
Annualised mean return & 33.66\% & 32.53\% & +1.13 pp \\
Annualised volatility & 30.17\% & 29.68\% & +0.49 pp \\
\bandrow{RISK-ADJUSTED}
Annualised Sharpe ratio & 1.12 & 1.10 & +0.02 \\
\end{metrictable}

\notesource{Source: skfolio Population summary, applied to out-of-sample returns.}

\notesection{Reproducibility}
\begin{notecode}
from skfolio import Population, RiskMeasure
from skfolio.optimization import InverseVolatility, RiskBudgeting
iv = InverseVolatility().fit(X_tr)
\end{notecode}

\notedivider

\begin{multicols}{2}
\glossaryitem{Drawdown.}{The percentage decline of cumulative wealth from a
previous high-water mark.}
\glossaryitem{Sharpe ratio.}{Annualised mean excess return divided by annualised
volatility.}
\end{multicols}

\begin{disclosures}
\textbf{Important disclosures.} The information contained in this communication
is not a research recommendation or ``investment research'' and is classified as
a Marketing Communication.

© 2026. All rights reserved. Document classification: General.
\end{disclosures}

\end{document}
